\documentclass[11pt]{article}
\usepackage[margin=1in]{geometry}
\usepackage{titlesec}
\usepackage{url}
\usepackage{hyperref}
\usepackage{amsmath}
\usepackage{enumitem} % in preamble
\setlist[itemize]{leftmargin=*, nosep, topsep=0pt, partopsep=0pt}

\newcommand{\memoheader}[4]{
  \noindent
  \textbf{From:} #2\\ 
  \textbf{To:} #1\\
  \textbf{Date:} #3\\
  \textbf{Subject:} #4\\[1ex]\hrule\vspace{1.5ex}
}

\begin{document}
\memoheader{Prof. Barbara Linke}
            {Miguel Cuaycong}
            {\today}
            {EME 150B: Lec 1}
\noindent This lecture was spent discussing course policies and the syllabus in general.
\\~\\
First homework due jan 16, 11:59pm
\section*{Design}
What is machine design?
\begin{itemize}
    \item Machine - System of elements arranged to transmit motion and energy in a predetermined and controlled fashion
    \item Design - Creative process to find a solution for a problem
    \item Design theory is a field that analyzes how we design rather than "what"("what" is what this class is about)
    \item 
\end{itemize}
\subsection*{Trends in product design}
\begin{itemize}
    \item Modern design is challenged by increasing complexity(number of parts and steps)
    \item Most complex products incorporate mechanical, electrical components. Also software, control modules, and human-machine interfaces
    \item miniaturization can also cause complexity
    \item Sustainability and repairability (some designs make it easier for consumers to do a DIY repair)
\end{itemize}
\subsection*{Design Process}
\begin{enumerate}
    \item General steps of design process:
    \begin{itemize}
        \item Identification of need(Market demand, noise, etc)
    \end{itemize}
    \item Definition of problem
    \item Concept design
    \item Analysis and optimization
    \item evaluation
    \item last three steps are iterative
\end{enumerate}
\subsection*{Factor of safety}
\subsection*{Cost and life cycle}
\section*{Gears and gear trains}
\begin{itemize}
    \item Purpose of gears:
    \\
    - Components that transmit energy through rotational force from one gear/device to another gear/device
    \item History of gears:
    \\
    - Water or wind powered wheels
    \\- Leonardo Da Vinci during 1500
    \item Mesing:
    \\
    -conjugate screw rotor
    \\-cycloid
    \item Types of gears
    \\-Internal gear: e.g. Planetary gear
    \\-External gear
    \\-Spur gear
    \\-Helical gear
    \\-Bevel gear
    \\-Worm gear
    \\-Rack and pinion: (Pinion is the smaller of one or two gear)
    \\-
    \\-
    \\-
    \\-
    \\-
\end{itemize}
\end{document}
